\chapter{Conclusão}
\label{ch:conclusao}

O Capítulo~\ref{ch:conclusao} apresenta as principais conclusões da
dissertação. A Secção~\ref{sec:sumario} sintetiza as contribuições,
revisita os objetivos enunciados em
Secção~\ref{sec:objetivos} e regista o estado de execução de cada
um. A Secção~\ref{sec:trabalho-futuro} identifica linhas de trabalho
futuro.

\noindent\textbf{Conceitos-chave:} conclusão; síntese;
contribuições; \emph{research questions} respondidas; trabalho
futuro.

% ---------------------------------------------------------------------------
\section{Sumário e Principais Contribuições}
\label{sec:sumario}

Esta dissertação investigou a viabilidade de aplicar algoritmos de
ranqueamento baseados em grafos --- concretamente PageRank, HITS e
SALSA --- à priorização estrutural de requisitos de software. O
problema original, exposto no Capítulo~\ref{ch:introducao}, partia da
constatação de que as técnicas tradicionais de Priorização de
Requisitos de Software (PRS), como a Votação Acumulativa, o AHP e o
PERT/CPM, são fortemente dependentes da intervenção subjetiva dos
\emph{stakeholders} e não escalam para sistemas com elevado número
de requisitos interdependentes
\parencite{silva2018, achimugu2014, bukhsh2020}.

A resposta proposta consistiu em transferir parte do esforço de
priorização para a topologia do sistema: se o grafo de dependências
contém informação latente sobre a importância relativa dos
requisitos, então um algoritmo de ranqueamento adequadamente
escolhido pode extrair essa informação de forma reprodutível e
automatizada. Para validar esta hipótese, foram desenvolvidos dois
artefactos:

\begin{itemize}
    \item \textbf{ReqGraph} --- um protótipo Python, organizado
        como pacote instalável, que automatiza o \emph{pipeline}
        \emph{Código-fonte} $\rightarrow$ AST $\rightarrow$
        \emph{Call Graph} $\rightarrow$ \emph{Grafo de Requisitos}.
        A análise é totalmente estática (não requer execução do
        código) e o protótipo gera artefactos em três formatos
        canónicos (PNG, JSON e DOT/Graphviz).
    \item \code{ranker.py} --- um módulo de análise que aplica
        PageRank, HITS e SALSA aos grafos produzidos pelo
        \emph{ReqGraph}, gerando \emph{scores} comparáveis, gráficos
        de barras e relatórios consolidados. A partir de R5, o
        \emph{script} \code{compute\_rank\_correlations.py} acrescenta
        ao \emph{pipeline} o cálculo de Spearman $\rho$ e Kendall
        $\tau_b$ entre cada par de algoritmos.
\end{itemize}

A validação foi realizada com \textbf{três projetos reais de código
aberto} (CPython \emph{stdlib}, Flask e \emph{scikit-learn}),
selecionados como representantes de regimes crescentes de
acoplamento arquitetural. Os principais resultados, detalhados no
Capítulo~\ref{ch:resultados}, sustentam quatro contribuições:

\begin{enumerate}
    \item \textbf{Validação da qualidade dos rankings extraídos
        automaticamente.} O \emph{pipeline} \emph{ReqGraph} produziu,
        sem intervenção humana adicional ao mapeamento
        \code{func\_to\_req}, grafos de requisitos coerentes com a
        arquitetura conhecida dos três projetos --- incluindo a
        deteção emergente do ciclo
        \req{REQUEST\_HANDLING}$\leftrightarrow$\req{CONTEXT}$\leftrightarrow$\req{ROUTING}$\leftrightarrow$\req{ERROR\_HANDLING}
        no Flask e da centralidade de \req{REQ\_PIPELINE} no
        \emph{scikit-learn}. A validação realizada incidiu sobre a
        \textbf{qualidade dos resultados} (coerência com a
        arquitetura conhecida), e \emph{não} sobre a sua adequação
        a perceções de profissionais --- limitação assumida e
        encaminhada para trabalho futuro
        (Secção~\ref{sec:trabalho-futuro}).

    \item \textbf{Comparação empírica quantitativa entre PageRank,
        HITS e SALSA} sobre grafos reais de requisitos, com
        métricas formais de \emph{rank correlation}. Os resultados
        (Tabela~\ref{tab:correlacoes}) mostram que a concordância
        entre algoritmos é \emph{parcial} e \emph{seletiva}: apenas
        dois dos $18$ pares analisados satisfazem o critério forte
        de convergência ($\rho \geq 0{,}7$ \emph{e} $\tau_b \geq 0{,}55$),
        ambos no caso mais esparso (\emph{stdlib}). Esta
        descoberta \emph{contradiz} a hipótese inicial de que
        densidade implicaria convergência --- e é, por si só,
        contribuição empírica relevante. A interpretação adequada é
        que cada algoritmo capta uma dimensão diferente de
        centralidade: PageRank capta importância transitiva, HITS
        \emph{authority} capta dependência direta, HITS \emph{hub}
        capta orquestração.

    \item \textbf{Implementação manual e documentada do SALSA.}
        Dado que o NetworkX não fornece SALSA, foi desenvolvida uma
        implementação completa baseada em matrizes de transição de
        cadeias de Markov, com tratamento explícito de vértices
        absorventes. Esta implementação é reaproveitável em
        contextos académicos e industriais.

    \item \textbf{Mecanismo de redução do custo de mapeamento.} O
        ficheiro \code{llm\_prompt\_mapping.md} formaliza um
        \emph{prompt} estruturado para gerar o dicionário
        \code{func\_to\_req} via modelos de linguagem de grande
        escala, mitigando uma das fricções identificadas --- a
        necessidade de mapear manualmente centenas de funções para
        domínios de requisito.
\end{enumerate}

\paragraph{Estado dos objetivos e \emph{research questions}.} A
Tabela~\ref{tab:estado} revisita os objetivos da
Secção~\ref{sec:objetivos} e identifica o capítulo de resposta a
cada uma das \emph{research questions}.

\begin{table}[ht]
    \centering
    \caption{Estado dos objetivos e das \emph{research questions}.}
    \label{tab:estado}
    \begin{tabular}{@{}l l l@{}}
        \toprule
        \textbf{Item} & \textbf{Estado} & \textbf{Capítulo / Secção} \\
        \midrule
        Objetivo (i) Formalização teórica
            & Concluído
            & Capítulo~\ref{ch:fundamentos} \\
        Objetivo (ii) Protótipo \emph{ReqGraph}
            & Concluído (\emph{open source})
            & Capítulo~\ref{ch:metodologia}, Secção~\ref{sec:reqgraph} \\
        Objetivo (iii) Avaliação empírica
            & Concluído sobre três casos
            & Capítulo~\ref{ch:resultados} \\[4pt]

        RQ1 (concordância entre algoritmos)
            & Respondida quantitativamente
            & Secção~\ref{sec:res-correlacoes} \\
        RQ2 (efeito da densidade)
            & Respondida (\emph{contraintuitiva})
            & Secção~\ref{sec:res-correlacoes}, \ref{sec:discussao} \\
        RQ3 (coerência arquitetural)
            & Respondida qualitativamente
            & Secção~\ref{sec:discussao} \\[4pt]

        Validação humana dos rankings
            & Não realizada (limitação)
            & Trabalho futuro (Secção~\ref{sec:trabalho-futuro}) \\
        Análise de vulnerabilidades topológicas
            & Apenas teórica
            & Capítulo~\ref{ch:fundamentos}, Secção~\ref{sec:salsa} \\
        Modelo híbrido (vértices artificiais)
            & Não implementado
            & Trabalho futuro \\
        \bottomrule
    \end{tabular}
\end{table}

A consistência entre as conclusões do referencial teórico e os
resultados experimentais sustenta a tese de que os algoritmos de
ranqueamento estrutural constituem um \textbf{complemento} --- não
um substituto --- às técnicas tradicionais de priorização,
contribuindo para a \textbf{redução, não a eliminação}, do viés
subjetivo em sistemas de larga escala. A subjetividade é deslocada,
do julgamento de prioridade por \emph{stakeholder} para a
construção do mapeamento \code{func\_to\_req}, com vantagens
operacionais que esta dissertação documenta empiricamente.

% ---------------------------------------------------------------------------
\section{Trabalho Futuro}
\label{sec:trabalho-futuro}

A investigação realizada abre várias linhas de continuação,
ordenadas por proximidade ao trabalho desenvolvido e endereçando
diretamente as ameaças à validade enunciadas em
Secção~\ref{sec:ameacas}.

\begin{itemize}
    \item \textbf{Validação humana dos rankings.} A presente
        avaliação é estrutural e qualitativa. Um estudo controlado
        com engenheiros de software --- produzindo rankings manuais
        sobre os mesmos projetos e comparando-os com os automáticos
        via correlação de Spearman ou Kendall --- forneceria uma
        medida quantitativa direta de adequação. Esta validação foi
        conduzida por \citet{silva2018} para o PageRank
        isoladamente; a sua replicação para o trio PageRank/HITS/SALSA,
        alargada a um corpus mais amplo, foi inviabilizada pela
        limitação de recursos disponível na presente investigação.

    \item \textbf{\emph{Inter-rater reliability} do mapeamento
        \code{func\_to\_req}.} Replicar o protocolo
        (Secção~\ref{sec:protocolo-func2req}) com um segundo
        investigador independente, reportando Cohen's $\kappa$ e
        analisando o impacto das divergências de mapeamento sobre
        os rankings finais. Esta extensão endereça diretamente a
        ameaça à validade interna principal desta dissertação.

    \item \textbf{Modelo Híbrido com Vértices Artificiais.}
        Implementar a técnica de injeção de prioridades de negócio
        através de vértices artificiais conectados ao grafo de
        requisitos \parencite{silva2018}. Esta extensão permitirá
        medir empiricamente o impacto de prioridades subjetivas
        sobre os rankings estruturais e oferecer um mecanismo de
        ajuste fino aos \emph{stakeholders}.

    \item \textbf{Análise de sensibilidade ao $d$ do PageRank.}
        Sistematizar a robustez dos rankings face a
        $d \in \{0{,}5,\, 0{,}7,\, 0{,}85,\, 0{,}95\}$ e às
        tolerâncias de convergência, produzindo gráficos de
        estabilidade de \emph{rank} por requisito.

    \item \textbf{Análise sobre \emph{codebase} completo.} Reaplicar
        o \emph{ReqGraph} aos três projetos sem seleção de
        ficheiros, reduzindo o número de \emph{singletons} e
        permitindo observar se os rankings dos requisitos
        atualmente isolados (\eg\ \req{REQ\_BLUEPRINTS},
        \req{REQ\_TEMPLATING} no Flask) alteram a coerência
        arquitetural geral.

    \item \textbf{Extensão a outras linguagens.} O motor de
        extração de \emph{Call Graph} baseia-se exclusivamente no
        módulo \code{ast} do Python. Estender o \emph{ReqGraph} a
        Java (via \code{JavaParser}), TypeScript (via TypeScript
        Compiler API) ou C/C++ (via \code{libclang}) permitiria
        validar a generalidade da abordagem para além do
        ecossistema Python.

    \item \textbf{Integração contínua e \emph{snapshots}
        temporais.} Aplicar o \emph{ReqGraph} a \emph{snapshots}
        sucessivos do mesmo projeto (\eg\ \emph{tags} de
        \emph{releases}) permitiria estudar a \textbf{evolução
        estrutural} dos requisitos ao longo do tempo, detetando
        \emph{drift} arquitetural e regressões de acoplamento.

    \item \textbf{Tratamento de pesos semânticos.} Atualmente, todas
        as arestas têm peso unitário. Atribuir pesos derivados da
        frequência de invocação ou da força do acoplamento
        (\eg\ número de funções partilhadas entre dois requisitos)
        poderá refinar substancialmente os rankings.
\end{itemize}

O conjunto destas linhas confirma que, embora o problema central
tenha sido endereçado, a abordagem proposta abre um espaço de
investigação largo no cruzamento entre Engenharia de Requisitos,
Teoria dos Grafos e Análise Estática de Código.
